\documentclass[10pt,a4paper]{article}
\usepackage[T1]{fontenc}
\usepackage[utf8]{inputenc}
\usepackage[brazil]{babel}
\AtBeginDocument{\shorthandoff{"}}
\usepackage[margin=2cm]{geometry}
\usepackage{amsmath,amssymb}
\usepackage{booktabs,tabularx,array,longtable}
\usepackage{tikz}
\usetikzlibrary{positioning,arrows.meta,shapes.misc,calc}
\usepackage{hyperref}
\newcolumntype{L}{>{\raggedright\arraybackslash}X}
\newcolumntype{P}[1]{>{\raggedright\arraybackslash}p{#1}}
\newcommand{\cod}[1]{{\ttfamily\small\def\_{\textunderscore\allowbreak}#1}}
\renewcommand{\arraystretch}{1.15}
\setlength{\parskip}{3pt}

\title{Anexo E1 — Caderno de Especificação de Design do UX-Agent\\\large Chat de estudo para estudantes universitários autistas (roteiro-0.2)}
\author{Matheus Emanuel da Silva \and João Francisco Teles da Silva \and Eduardo Almeida\\CEFET-MG, Campus Divinópolis — Interação Humano-Computador}
\date{Novembro de 2026}

\begin{document}
\maketitle

\noindent Este caderno especifica o Projeto de Interação (Parte I) que a implementação (Parte II) realiza. Cada artefato cita o identificador usado no código e no log, para que a rastreabilidade entre design, código e registros possa ser conferida. A persona, as tarefas C1, C2, V1 e V2 e a HTA vêm da modelagem congelada; V1 e V2 aparecem aqui só como especificação, nunca foram executadas.

\tableofcontents

\section{Metamensagem}
A interface realiza a metamensagem do designer para dois receptores ao mesmo tempo: a pessoa estudante e o agente proxy que a substitui no piloto.

\begin{quote}\itshape
``Entendo que você é estudante universitária, usa web e chat com frequência, estuda sozinha e quer saber exatamente o que, quando e em que formato entregar, sem ter de perguntar a um colega e sem ser tratada como incapaz. Entendo também que respostas de chat podem trazer metáfora, ironia ou instrução implícita, e que você pode querer conferir o sentido antes de agir. Por isso, este chat mostra a tarefa num cartão fixo, guarda tudo no histórico, avisa quando a mensagem continua, deixa você pedir uma explicação em lista quando quiser, oferece ajuda sobre como fazer e como corrigir, aceita correções de campo sem punição e pergunta antes de encerrar com campo vazio. Você decide quando terminou; quem confere o resultado é um verificador externo, não o próprio chat.''
\end{quote}

Para o agente, a mesma metamensagem é entregue pela observação $o_t$: ele recebe os mesmos signos, na mesma ordem, e nenhum signo a mais. A tabela de paridade (Seção~\ref{sec:paridade}) mostra como cada signo chega a cada receptor.

\section{Tabela de paridade de signos}\label{sec:paridade}
Tipos de signo segundo a Engenharia Semiótica: \emph{estático} (não muda durante a cena), \emph{dinâmico} (muda com a interação) e \emph{metalinguístico} (fala sobre a própria interface). O tradutor é \cod{uxagent/observacao.py}; a tela humana é \cod{uxagent/web.py}; ambos leem a mesma função \cod{EstadoInterface.tela()}.

\begin{longtable}{@{}P{2.6cm}P{1.6cm}P{3.6cm}P{4.2cm}P{3.6cm}@{}}
\toprule
Signo & Tipo & Tela humana & Observação $o_t$ (agente) & Vazamento evitado\\
\midrule\endhead
Título ``Chat de estudo'' & estático & cabeçalho & \cod{TELA: Chat de estudo} & --\\
Cartão da tarefa & estático & caixa fixa no topo & \cod{CARTÃO DA TAREFA: \ldots} (texto idêntico) & instrução escrita a partir do objetivo, sem o caminho da solução\\
Rótulos dos campos & estático & rótulo de cada campo & nomes em \cod{CAMPOS DE RESPOSTA} e no \cod{enum} de \cod{campo} & só os campos da tarefa; sem formato esperado\\
Valor dos campos & dinâmico & caixa de texto preenchida & \cod{Local: "\ldots"} ou \cod{(vazio)} & sem marca de certo/errado\\
Mensagens do chat & dinâmico & balões no histórico & \cod{[Chat] \ldots} na mesma ordem & sem rótulo de figura (metáfora, ironia, implícita)\\
Indicador de continuação & dinâmico & ``A mensagem continua abaixo'' & mesma frase, só quando há bloco oculto & não revela o conteúdo do bloco B\\
Mensagens do sistema & dinâmico & linha cinza (preenchido, substituído) e aviso (recusado, confirmação) & \cod{[Sistema] \ldots} e \cod{AVISO DO SISTEMA} & não avalia o conteúdo enviado\\
Contador de ações & dinâmico & ``Ações restantes: n de 12'' & \cod{AÇÕES RESTANTES: n de 12} & --\\
Alerta de poucas ações (AL) & dinâmico & contador em destaque & ``(atenção: restam poucas ações)'' & --\\
Botões & estático por cena & botões visíveis na cena & lista \cod{AÇÕES DISPONÍVEIS AGORA} com id e rótulo & só as falas $u{:}$ habilitadas pela MoLIC\\
Descrição dos botões & metalinguístico & rótulo do botão & frase curta após o rótulo & não diz quando usar\\
Texto de ajuda & metalinguístico & painel de ajuda & \cod{JANELA DE AJUDA ABERTA} (texto idêntico) & explica a interface, não a tarefa\\
Reformulação & dinâmico & nova mensagem do chat em lista & \cod{[Chat] \ldots} (idêntica) & só aparece depois de \cod{reformular}\\
\midrule
\multicolumn{5}{@{}l}{\emph{Signos e dados que não existem em nenhum dos receptores}}\\
\midrule
Rótulo de figura, texto superficial & -- & ausente & ausente; usado só pelo controlador & R1 lê \cod{figura} no roteiro, fora do prompt\\
Chave de sucesso, resposta canônica & -- & ausente & ausente & em \cod{analista/gabaritos.json}, nunca importado por \cod{uxagent/agente.py}\\
HTA, CTT, KLM, caminho mínimo & -- & ausente & ausente & instrumentos do analista\\
Parâmetros sorteados ($\theta$) & -- & ausente & ausente & gravados no log, não no prompt\\
\bottomrule
\end{longtable}

\paragraph{Desvios declarados.} Dois mecanismos de perfil alteram só $o_t$ e são registrados como \cod{restricao\_aplicada}: (i) o filtro de leitura (R1) troca o texto integral de um bloco com figura pelo superficial; (ii) os filtros de \emph{affordance} (R2, R3, R4 e limite de tentativas) retiram falas da lista por um passo. A pessoa sempre vê a tela integral. O agente recebe ainda a lista das próprias ações anteriores na tarefa, equivalente à memória de curto prazo que a pessoa teria; o contexto é reiniciado entre tarefas.

\section{Refinamento das tarefas em CTT}
Operadores (Dix et al.): $\gg$ habilitação; $[]$ escolha; $|\!=\!|$ independência de ordem; $[\!>$ desativação; $|\!>$ suspensão e retomada; $^{*}$ iteração; $[T]$ opcional. Tipos: abstrata (A), do usuário/cognitiva (U), de interação (I), da aplicação (S). A CTT é leitura do analista e não é fornecida ao agente.

\subsection{Estrutura comum}
\begin{align*}
\text{Tarefa}_X &:= \big(\text{Entender} \gg \text{ObterInfo} \gg \text{Registrar} \gg \text{Concluir}\big)\ [\!>\ \text{Interrupção}_S\\
\text{Tarefa}_X &\ |\!>\ \text{Ajuda}\\
\text{Ajuda} &:= (\text{AbrirComoFazer}_I\ []\ \text{AbrirComoCorrigir}_I) \gg \text{LerAjuda}_U \gg \text{Voltar}_I\\
\text{Entender} &:= \text{LerCartão}_U \gg \text{DecidirOQueProcurar}_U\\
\text{Concluir} &:= \text{Conferir}_U \gg \text{Encerrar}_I \gg [\text{Confirmação}]\\
\text{Confirmação} &:= \text{PedeConfirmação}_S \gg (\text{Confirmar}_I\ []\ \text{Cancelar}_I)\\
\text{Interrupção}_S &:= \text{limite de 12 ações}\ []\ \text{falha técnica}
\end{align*}

\subsection{Por tarefa}
\begin{tabularx}{\textwidth}{@{}P{1.2cm}LP{1.6cm}@{}}
\toprule
Tarefa & ObterInfo e Registrar & Conjunto\\
\midrule
C1 & ObterInfo $:=$ LerFala$_U$ $\gg$ ([NovaPergunta$_I$] $[]$ [Reformular$_I$])$^{*}$ \newline Registrar $:=$ EnviarLocal$_I$ $|\!=\!|$ EnviarHorário$_I$ & calibração\\
C2 & ObterInfo $:=$ LerBlocoA$_U$ $\gg$ Rolar$_I$ $\gg$ LerBlocoB$_U$ $\gg$ InterpretarFiguras$_U$ $\gg$ ([Reformular$_I$] $[]$ [NovaPergunta$_I$])$^{*}$ \newline Registrar $:=$ EnviarEntrega$_I$ $|\!=\!|$ EnviarPrazo$_I$ & calibração\\
V1 & ObterInfo $:=$ LerBlocoA$_U$ $\gg$ Rolar$_I$ $\gg$ LerBlocoB$_U$ $\gg$ ([Reformular$_I$] $[]$ [NovaPergunta$_I$])$^{*}$ \newline Registrar $:=$ EnviarLocal$_I$ $|\!=\!|$ EnviarPrazo$_I$ & validação (congelada)\\
V2 & ObterInfo $:=$ LerBlocoA$_U$ $\gg$ Rolar$_I$ $\gg$ LerBlocoB$_U$ $\gg$ InferirPrazo$_U$ $\gg$ ([Reformular$_I$] $[]$ [NovaPergunta$_I$])$^{*}$ \newline Registrar $:=$ EnviarPrazo$_I$ $|\!=\!|$ EnviarFormato$_I$ & validação (congelada)\\
\bottomrule
\end{tabularx}

\medskip
\begin{center}
\begin{tikzpicture}[font=\scriptsize, every node/.style={draw, rounded corners=2pt, align=center, inner sep=2pt}, level distance=14mm,
  level 1/.style={sibling distance=43mm}, level 2/.style={sibling distance=15mm}, edge from parent/.style={draw,-}]
\node (r) {Tarefa C2 (A)}
  child {node (e) {Entender (A)}
    child {node {Ler\\cartão\\(U)}}
    child {node {Decidir\\(U)}}}
  child {node (o) {ObterInfo (A)}
    child {node {Ler A\\(U)}}
    child {node {Rolar\\(I)}}
    child {node {[Reform.]*\\(I)}}}
  child {node (g) {Registrar (A)}
    child {node {Entrega\\(I)}}
    child {node {Prazo\\(I)}}}
  child {node (c) {Concluir (A)}
    child {node {Conferir\\(U)}}
    child {node {Encerrar\\(I)}}};
\path (e) -- node[draw=none, fill=white]{$\gg$} (o);
\path (o) -- node[draw=none, fill=white]{$\gg$} (g);
\path (g) -- node[draw=none, fill=white]{$\gg$} (c);
\path (g-1) -- node[draw=none, fill=white]{$|\!=\!|$} (g-2);
\node[draw=none] at (0,-4.1) {Raiz: $(\ldots)\ [\!>\ \text{Interrupção}_S$ \quad e \quad $(\ldots)\ |\!>\ \text{Ajuda}$};
\end{tikzpicture}
\end{center}

\section{Linha de base GOMS/KLM de especialista}
Cena principal (\cod{conversa}), método do especialista sem ajuda, igual ao caminho mínimo da HTA. Operadores e tempos de Card, Moran e Newell: $K=0{,}28$\,s (digitador médio), $P=1{,}10$\,s, $B=0{,}10$\,s, $H=0{,}40$\,s, $M=1{,}35$\,s. Regra de colocação de $M$: um $M$ antes de cada fala $u{:}$. Clique em botão $= M\,P\,B$; envio de campo $= M\,P\,B\,H\,K^{n}\,H\,P\,B$, com $n$ igual ao comprimento da resposta canônica. Leitura do texto não tem operador no KLM e fica fora da linha de base (ameaça declarada no artigo). Gerado por \cod{python -m analista.klm}.

\begin{tabularx}{\textwidth}{@{}P{1cm}P{1.2cm}LP{1.5cm}@{}}
\toprule
Tarefa & $A_{\min}$ & Sequência de operadores & $T_{\text{KLM}}$\\
\midrule
C1 & 3 & enviar Local ($MPBHK^{13}HPB$ = 8,19\,s) $\gg$ enviar Horário ($MPBHK^{9}HPB$ = 7,07\,s) $\gg$ encerrar ($MPB$ = 2,55\,s) & 17,81\,s\\
C2 & 4 & rolar ($MPB$ = 2,55) $\gg$ enviar Entrega ($MPBHK^{38}HPB$ = 15,19) $\gg$ enviar Prazo ($K^{9}$, 7,07) $\gg$ encerrar (2,55) & 27,36\,s\\
V1 & 4 & rolar (2,55) $\gg$ enviar Local ($K^{16}$, 9,03) $\gg$ enviar Prazo ($K^{9}$, 7,07) $\gg$ encerrar (2,55) & 21,20\,s\\
V2 & 4 & rolar (2,55) $\gg$ enviar Prazo ($K^{12}$, 7,91) $\gg$ enviar Formato ($K^{9}$, 7,07) $\gg$ encerrar (2,55) & 20,08\,s\\
\bottomrule
\end{tabularx}

\medskip
O tempo simulado do agente usa a mesma regra de operadores para cada fala executada e soma a hesitação sorteada pelo controlador (\cod{uxagent/tempo.py}). A latência da API do LLM é gravada à parte (\cod{latencia\_llm\_s}) e nunca entra no tempo simulado.

\section{Diagrama MoLIC}
Papel de usuário único: estudante (representado pela persona E. ou por uma pessoa). Cenas em retângulos arredondados; processamento do sistema em caixas pretas; falas $u{:}$ do usuário e $d{:}$ do preposto nas transições; a ajuda é acesso ubíquo a partir da cena principal. A caixa preta tracejada é a verificação do limite de ações feita pelo sistema após cada fala $u{:}$.

\begin{center}
\begin{tikzpicture}[font=\scriptsize, >=Stealth,
  cena/.style={draw, rounded corners=6pt, align=center, minimum width=36mm, minimum height=15mm, inner sep=3pt},
  proc/.style={fill=black, minimum width=6mm, minimum height=4mm, inner sep=0},
  ponto/.style={circle, fill=black, minimum size=3mm, inner sep=0},
  fim/.style={circle, draw, double, fill=black, minimum size=3mm, inner sep=0}]
\node[ponto] (ini) at (0,4.3) {};
\node[cena] (conv) at (0,0) {\textbf{Conversa}\\Tópico: completar a tarefa\\ler cartão; ler histórico;\\preencher campos};
\node[cena] (aj) at (-7.2,0) {\textbf{Ajuda}\\Tópico: tirar dúvida\\como fazer; como corrigir};
\node[proc] (p1) at (5.0,2.2) {};
\node[proc] (p2) at (5.0,-2.4) {};
\node[cena] (conf) at (5.0,-5.6) {\textbf{Confirmação}\\Tópico: encerrar\\com campo vazio?};
\node[fim] (fim) at (8.0,-2.4) {};
\node[proc] (p3) at (-3.6,-3.6) {};
\draw[->] (ini) -- node[right]{d: mostra cartão e 1ª mensagem} (conv);
\draw[->] (conv.170) -- node[above, align=center]{u: ajuda\_como\_fazer\\u: ajuda\_corrigir} (aj.10);
\draw[->] (aj.-10) -- node[below]{u: voltar / d: retoma} (conv.190);
\draw[->] (conv.20) -- node[above left]{u: enviar} (p1);
\draw[->] (p1) |- node[pos=0.25, right, align=left]{d: campo preenchido\\d: substituído (RA)\\d: recusado (CE)} (conv.5);
\draw[->] (conv.-100) .. controls +(-0.6,-2.2) and +(-2.4,-0.6) .. node[pos=0.55, below left, align=right]{u: rolar\_historico / d: revela bloco\\u: reformular / d: lista literal\\u: nova\_pergunta / d: fala neutra} (conv.-165);
\draw[->] (conv.-20) -- node[below left]{u: encerrar} (p2);
\draw[->] (p2) -- node[right]{d: há campo vazio (PA)} (conf);
\draw[->] (p2) -- node[above]{d: encerra} (fim);
\draw[->] (conf.east) -| node[pos=0.75, left, align=right]{u: confirmar\_encerrar\\d: encerra} (fim);
\draw[->] (conf.west) -| node[pos=0.25, below]{u: cancelar\_encerrar} (1.2,-0.75);
\draw[->, dashed] (conv.-130) -- (p3);
\draw[->, dashed] (p3) -- node[below, align=center]{12 ações ou falha técnica\\d: interrompe} (2.0,-3.6) -- (fim.south west);
\end{tikzpicture}
\end{center}

\begin{longtable}{@{}P{2.7cm}P{1.8cm}P{3.6cm}P{3.2cm}P{4.2cm}@{}}
\toprule
Fala $u{:}$ (id do agente) & Cena & Fala $d{:}$ (id no log) & Evento no log & Mudança de estado da interface\\
\midrule\endhead
\cod{enviar}(campo, texto) & conversa & \cod{d\_campo\_preenchido}, \cod{d\_campo\_substituido}, \cod{d\_texto\_vazio\_recusado}, \cod{d\_campo\_inexistente} & \cod{acao\_executada}; \cod{erro\_capturado} & valor do campo; linha no histórico\\
\cod{rolar\_historico} & conversa & \cod{d\_revela\_bloco} & \cod{acao\_executada} & bloco B visível; some o indicador\\
\cod{reformular} & conversa & \cod{d\_reformulacao} & \cod{acao\_executada}, \cod{ajuda\_solicitada} (significado) & lista literal no histórico\\
\cod{nova\_pergunta} & conversa & \cod{d\_fala\_neutra} & \cod{acao\_executada} & fala neutra no histórico\\
\cod{ajuda\_como\_fazer} & conversa & \cod{d\_abre\_ajuda} & \cod{ajuda\_solicitada} (procedimento) & cena ajuda\\
\cod{ajuda\_corrigir} & conversa & \cod{d\_abre\_ajuda} & \cod{ajuda\_solicitada} (recuperação) & cena ajuda\\
\cod{voltar} & ajuda & \cod{d\_fecha\_ajuda} & \cod{acao\_executada} & volta à conversa no mesmo estado\\
\cod{encerrar} & conversa & \cod{d\_pede\_confirmacao} ou \cod{d\_encerra} & \cod{acao\_executada}; \cod{encerrada} & cena confirmação ou fim\\
\cod{confirmar\_encerrar} & confirmação & \cod{d\_encerra} & \cod{encerrada} & fim com campo vazio\\
\cod{cancelar\_encerrar} & confirmação & \cod{d\_fecha\_ajuda} & \cod{acao\_executada} & volta à conversa\\
(sistema) & qualquer & \cod{d\_alerta\_acoes}, \cod{d\_interrompe} & \cod{alerta\_exibido}, \cod{interrompida} & destaque no contador; fim\\
\bottomrule
\end{longtable}

\noindent Implementação: \cod{uxagent/molic.py} define as cenas e \cod{falas\_habilitadas(estado)}; o esquema JSON da resposta do LLM usa essa lista como \cod{enum}, então o agente não consegue escolher uma fala fora da cena.

\section{Matriz de prevenção e recuperação de erros por operação da HTA}
PP: prevenção passiva (o signo evita o erro); PA: prevenção ativa (o sistema impede ou pede confirmação); AL: prevenção apoiada (alerta antes do erro); RA: recuperação apoiada (o sistema ajuda a desfazer); CE: captura de erro (o sistema detecta e recusa).

\begin{longtable}{@{}P{2.1cm}P{2.6cm}P{2.6cm}P{2.3cm}P{2.6cm}P{2.6cm}@{}}
\toprule
Operação HTA & PP & PA & AL & RA & CE\\
\midrule\endhead
1.1 Ler cartão & cartão fixo no topo de toda cena & -- & -- & -- & --\\
2.1 Ler bloco A & histórico persistente & -- & -- & reformular repete em lista & --\\
2.2 Rolar histórico & indicador ``a mensagem continua''; botão só existe com bloco oculto & -- & -- & -- & --\\
2.3 Reformular & rótulo ``Explicar/reformular'' & -- & -- & a lista substitui a dúvida & --\\
2.4 Nova pergunta & -- & -- & -- & fala neutra não apaga nada & --\\
3.1/4.1 Enviar campo & rótulos dos campos; \cod{enum} de campos & campo inexistente não é aceito & -- & substituir valor mostra o anterior; ajuda ``como corrijo?'' & texto vazio recusado (\cod{erro\_capturado})\\
5.2 Encerrar & contador visível & confirmação se há campo vazio (\cod{d\_pede\_confirmacao}) & alerta com 3 ações restantes (\cod{alerta\_exibido}) & ``Voltar e completar'' & --\\
Qualquer (agente) & lista de falas habilitadas & proposta fora da lista rejeitada & -- & -- & \cod{acao\_rejeitada}; 3 seguidas: \cod{falha\_tecnica}\\
\bottomrule
\end{longtable}

\noindent Assinaturas no log usadas por \cod{analista/analisar.py}: PA = \cod{acao\_executada} com \cod{d\_pede\_confirmacao}; AL = \cod{alerta\_exibido}; RA = \cod{d\_campo\_substituido} ou \cod{ajuda\_corrigir}; CE = \cod{erro\_capturado}; PP = passo em que \cod{rolar\_historico} não estava habilitado.

\section{Ajuda sob demanda estratificada}
A ajuda só aparece depois de uma fala $u{:}$ registrada (\cod{ajuda\_solicitada}); nunca é empurrada ao agente. Os três estratos seguem os tipos de dúvida de Shneiderman e Plaisant.

\begin{tabularx}{\textwidth}{@{}P{2.4cm}P{3cm}P{2.6cm}LP{2.6cm}@{}}
\toprule
Tipo de dúvida & Pergunta da pessoa & Fala $u{:}$ & Conteúdo & O que não revela\\
\midrule
Significado & ``O que isto quer dizer?'' & \cod{reformular} & última mensagem do chat em lista literal, roteirizada por tarefa & nada além do que a mensagem já dizia\\
Procedimento & ``Como faço?'' & \cod{ajuda\_como\_fazer} & como usar cartão, campos, Enviar, Rolar, Explicar e Encerrar & a resposta da tarefa, o campo certo\\
Recuperação & ``Como corrijo?'' & \cod{ajuda\_corrigir} & enviar novo texto para o mesmo campo substitui o anterior & se o valor atual está certo\\
\bottomrule
\end{tabularx}

\section{Rupturas de comunicação (MAC) e suas assinaturas}
As rupturas são etiquetadas só por padrões de eventos; o texto livre do LLM nunca é lido para isso.

\begin{tabularx}{\textwidth}{@{}P{2.6cm}LP{3.2cm}@{}}
\toprule
Etiqueta & Assinatura no log & Operação HTA\\
\midrule
Cadê? & \cod{enviar} antes de \cod{rolar\_historico} numa tarefa com bloco oculto & 2.2, 3.1\\
O que é isto? & \cod{ajuda\_solicitada} do tipo significado (\cod{reformular}) & 2.3\\
Socorro! & \cod{ajuda\_solicitada} do tipo procedimento & qualquer\\
Epa! & \cod{d\_campo\_substituido} ou \cod{ajuda\_corrigir} & 3.1/4.1\\
Assim não dá & \cod{nova\_pergunta} logo após \cod{reformular} & 2.4\\
E agora? & duas ações seguidas sem mudança de campo na cena conversa & plano 2\\
Para mim está bom & \cod{encerrada} com campo preenchido fora da chave (engano de avaliação) & 5.2\\
Desisto & \cod{restricao\_aplicada} do limite de tentativas, \cod{confirmar\_encerrar} ou \cod{interrompida} & 5.2\\
\bottomrule
\end{tabularx}

\section{Rastreabilidade com o código}
\begin{tabularx}{\textwidth}{@{}P{4.2cm}L@{}}
\toprule
Artefato & Arquivo\\
\midrule
Roteiro das tarefas (tela) & \cod{roteiros/calibracao.json}, \cod{roteiros/validacao.json} (bloqueado sem \cod{--validacao-congelada})\\
MoLIC e falas & \cod{uxagent/molic.py}\\
Estado da interface & \cod{uxagent/interface.py}\\
Tradutor de observação & \cod{uxagent/observacao.py}\\
Mecanismos e amostrador & \cod{uxagent/mecanismos.py}, \cod{config/theta.json}\\
Agente LLM & \cod{uxagent/agente.py}\\
Registrador & \cod{uxagent/registrador.py}\\
Verificador, KLM, MAC & \cod{analista/verificador.py}, \cod{analista/klm.py}, \cod{analista/analisar.py}, \cod{analista/gabaritos.json}\\
\bottomrule
\end{tabularx}
\end{document}
